\clearpage\fancyhead[L]{\small GAUSS AN BODE.}\hypertarget{native-356-lower}{}\label{native-356-lower}
\begin{center}\textbf{GAUSS AN BODE.}\par\vspace{4pt}\hrule\vspace{4pt}\fontsize{10}{12}\selectfont Astronomisches Jahrbuch für 1815. Seite 190\qquad194. Berlin 1812.\par\vspace{4pt}\hrule\vspace{4pt}\end{center}

\begin{flushright}Göttingen, 5. Mai 1812.\end{flushright}
\par\indent Ich möchte von Ihnen, wie im vergangenen Jahr, auch in diesem Jahr einige Beobachtungen der Pallas um die Zeit der Opposition (10. Juni) erhalten. Deshalb mache ich Sie besonders darauf aufmerksam und teile Ihnen mit, was ich bisher dazu beobachtet habe. Am 4. April habe ich diesen Planeten zum ersten Mal wieder gesehen, aber erst am 9. beobachtet. Die gestrige Beobachtung war ebenso wie die vom 9. April für die Bestimmung der Deklination ungünstig, daher lasse ich diese ganz weg. Der Planet ist kaum heller als ein Stern 10. Größe.

\begin{center}
\fontsize{10}{12}\selectfont\setlength{\tabcolsep}{3pt}\renewcommand{\arraystretch}{1.15}\begin{tabular}{c|c|c|c}
1812 & Mittlere Zeit & Rektasz. $\PallasSymbol$ & Nördl. Dekl. \\
\hline
April\ \ 9 & $11^{\mathrm u}\;32^{\mathrm m}\;27^{\mathrm s}$ & $169^\circ\;5'\;58''4$ & \\
Mai\ \ 2 & $10\;36\;6$ & $268\;42\;5.3$ & $21^\circ\;4'\;37''7$ \\
4 & $10\;24\;16$ & $268\;32\;27.4$ & \\
\end{tabular}
\end{center}
\SourceLeaf{357}\par\indent Herr \textsc{Nicolai} findet beim Vergleich mit den Elementen, nach denen er seine Ephemeride berechnet hat, folgende Unterschiede:

\begin{center}
\fontsize{10}{12}\selectfont\setlength{\tabcolsep}{3pt}\renewcommand{\arraystretch}{1.15}\begin{tabular}{c|c|c}
 & Rektasz. & Dekl. \\
\hline
April\ \ 9 & $+30''6$ & \\
Mai\ \ 2 & $+57.6$ & $-5''4$ \\
4 & $+49.4$ & \\
\end{tabular}
\end{center}

\par\indent Ich halte meine Beobachtungen für gut. Verglichen wurde mit 93, 96, 95 im Herkules nach \textsc{Piazzi}s Katalog. Ich neige dazu, den Unterschied in der Rektaszension den Störungen durch Saturn und Mars zuzuschreiben. Sie sind in meinen Rechnungen bisher unberücksichtigt geblieben, doch ich werde sie bald ebenfalls berechnen. Die Zahl aller von mir berechneten periodischen Störungen durch Jupiter beläuft sich auf 400, wobei alle unter $1''$ weggelassen sind. Dabei haben sich mehrere bemerkenswerte Ergebnisse gezeigt, die ich aber vorerst noch zurückhalten muss. Denn meine ganze Rechnung kann nur als vorläufig gelten und muss noch einmal in größerem Umfang wiederholt werden.

\par\indent Von den anderen neuen Planeten habe ich noch keinen beobachtet. Ich habe gestern Abend nur die Gegend, in der Juno steht, ins Fernrohr genommen. Sie stand sehr tief, und die kleinen Sterne waren noch schlecht zu sehen. Es schien mir, als sei die Stelle, an der Juno nach der Ephemeride stehen sollte, leer. Bei einer um etwa $25'$ kleineren Rektaszension stand dagegen ein Stern 9. bis 10. Größe, der vielleicht Juno war. Das sei nur als vorläufige Notiz mitgeteilt. Bald wird sich das genauer bestimmen lassen.

\par\indent Bei dieser Gelegenheit teile ich Ihnen noch einige andere astronomische Beobachtungen mit. — Vom Erscheinen des zweiten Kometen erfuhr ich am 9. Dezember des vergangenen Jahres durch Herrn von \textsc{Lindenau}, und ich fand ihn gleich am selben Tag. Hier sind meine sämtlichen Beobachtungen dieses Kometen. Ich halte sie alle für gut, mit Ausnahme der letzten Deklination, die etwas weniger zuverlässig ist.

\begin{center}
\fontsize{10}{12}\selectfont\setlength{\tabcolsep}{3pt}\renewcommand{\arraystretch}{1.15}\begin{tabular}{c|c|c|c}
 & Mittlere Zeit & Rektasz. & Dekl. \\
\hline
1811 Dez.\ \ 9 & $10^{\mathrm u}\;6^{\mathrm m}\;52^{\mathrm s}$ & $63^\circ\;49'\;41''4$ & $10^\circ\;21'\;55''5$ S. \\
11 & $10\;34\;1$ & $63\;33\;18.0$ & $8\;39\;46.4$ \\
12 & $\;8\;5\;52$ & $63\;26\;25.3$ & $7\;54\;25.9$ \\
1812 Jan.\ \ 3 & $\;6\;54\;46$ & $62\;37\;6.4$ & $9\;47\;14.7$ N. \\
4 & $\;6\;54\;15$ & $62\;40\;38.6$ & $10\;29\;57.8$ \\
Febr.\ \ 2 & $11\;51\;52$ & $68\;9\;17.9$ & $26\;1\;34$ \\
\end{tabular}
\end{center}

\Needspace{15\baselineskip}\par\indent Gleich nach meiner Beobachtung vom 12. Dezember berechnete ich die parabolischen Elemente. Ich lasse sie hier jedoch weg und gebe stattdessen die von Herrn \textsc{Nicolai} verbesserten Elemente an, die dieser sehr geschickte Rechner gleich nach meiner Beobachtung vom 4. Januar ermittelt hat.

\begin{center}\fontsize{10}{12}\selectfont\setlength{\tabcolsep}{2pt}\renewcommand{\arraystretch}{1.15}\begin{tabular}{lll}
\multicolumn{3}{l}{Zeit des Perihels Nov 11.39211, Göttinger Meridian}\\
— Länge & $.\;.\;.\;.\;.\;.\;.\;.\;.\;.\;.\;.\;.\;.\;.$ & $47^\circ\;39''\;36''1$\\
— Logarithmus des Abstands & $.\;.\;.\;.\;.\;.\;.\;.$ & $0.2011781$\\
Länge des $\AscendingNode$ & $.\;.\;.\;.\;.\;.\;.\;.\;.\;.\;.\;.\;.\;.\;.\;.$ & $92^\circ\;53'\;44''2$\\
Bahnneigung & $.\;.\;.\;.\;.\;.\;.\;.\;.\;.\;.\;.\;.\;.$ & $31\;32\;38.7$\\
\multicolumn{3}{l}{Rechtläufige Bewegung.}\\
\end{tabular}\end{center}

\Needspace{10\baselineskip}\par\indent Der Vergleich dieser Elemente mit meinen Beobachtungen ergab nach Herrn \textsc{Nicolai}s Rechnung folgende Ergebnisse.
\SourceLeaf{358}\begin{center}
Differenz
\end{center}

\begin{center}
\fontsize{10}{12}\selectfont\setlength{\tabcolsep}{3pt}\renewcommand{\arraystretch}{1.15}\begin{tabular}{c|c|c||c|c|c}
1811 & Rektasz. & Dekl. & 1812 & Rektasz. & Dekl. \\
\hline
Dez. 9 & $-5''0$ & $+3''4$ & Jan. 3 & $-6''5$ & $+0''3$ \\
11 & $+2.9$ & $-32.7$ & 4 & $+2.8$ & $-12.7$ \\
12 & $+8.5$ & $-27.6$ & Febr. 2 & $-113.7$ & $+18.1$ \\
\end{tabular}
\end{center}

\par\indent Herr \textsc{Nicolai} will seine Elemente demnächst anhand der letzten im Februar gemachten Beobachtungen noch einmal verbessern.

\Needspace{30\baselineskip}\par\indent Hier noch einige Beobachtungen von Sternbedeckungen:

\begin{center}\fontsize{9}{11}\selectfont\setlength{\tabcolsep}{2pt}\renewcommand{\arraystretch}{1.15}\begin{tabular}{llrrrll}
1811 Okt. 23 & Stern 6. Größe & $7^{\mathrm u}$ & $21^{\mathrm m}$ & $40^{\mathrm s}1$ & Eintritt & Harding\\
— & 7. & 8 & 14 & 40.4 & — & Harding\\
— & 6.\ 7. & 8 & 30 & 33.8 & — & Gauss\\
24 & 7.\ 8. & 7 & 38 & 48.8 & — & Harding\\
1812 Febr. 1 & $\eta$ Virginis & 12 & 21 & 26.7 & — & Gauss\\
19 & $\gamma$ Tauri & 7 & 22 & 15.0 & Austritt & Gerling\\
— & 71 Tauri & 9 & 56 & 3.4 & Eintritt & Gerling\\
 &  &  &  & 3.8 & — & Harding\\
 &  &  &  & 3.9 & — & Gauss\\
— & $1\,\theta$ Tauri & 10 & 58 & 56.9 & — & Harding\\
 &  &  &  & 57.4 & — & Gauss\\
 &  &  &  & 57.9 & — & Gerling\\
 &  & 11 & 55 & 34.6 & Austritt & Gauss\\
 &  &  &  & 38.6 & — & Harding\\
— & $2\,\theta$ Tauri & 10 & 58 & 12.0 & Eintritt & Gauss und Harding\\
 &  & 11 & 59 & 6.6 & Austritt & Gauss\\
 &  &  &  & 8.6 & Austritt & Harding\\
— & 160 Mayer & 12 & 3 & 39.5 & Eintritt & Harding\\
 &  &  &  & 40.0 & — & Gauss\\
— & 162 Mayer & 12 & 26 & 27.0 & — & Gauss\\
 &  &  &  & 27.5 & — & Harding\\
20 & 111 Tauri & 11 & 7 & 46.8 & Eintritt & Gerling\\
— & Anonyma & 11 & 36 & 54.0 & — & Harding\\
— & 117 Tauri & 13 & 0 & 25.0 & — & Gerling\\
 &  &  &  & 25.1 & — & Harding\\
\end{tabular}\end{center}

\par\indent Die Bedeckung von $\eta\;\text{\AstroSymbols ♍}$ am 23. April fand hier in Göttingen gar nicht statt.

\par\indent Ich habe noch eine Beobachtung der Bedeckung von $\lambda\;\text{\AstroSymbols ♒}$ vom 2. September 1811 nachzutragen. Herr Oberst von \textsc{Tischleder} hat sie mir mit dem Auftrag zugesandt, sie Ihnen mitzuteilen.

\par\indent Festung Rosenberg oberhalb von Cronach: Eintritt von $\lambda$ Aquarii, 1811 Sept. 2, $10^{\mathrm u}$ $25^{\mathrm m}$ $5^{\mathrm s}4$ mittlere Zeit.

\par\indent Herr \textsc{Wachter} hat daraus und aus den Beobachtungen in Göttingen und Mannheim folgende Konjunktionszeiten berechnet: Mannheim $10^{\mathrm u}$ $16^{\mathrm m}$ $38^{\mathrm s}01$, Göttingen $10^{\mathrm u}$ $22^{\mathrm m}$ $29^{\mathrm s}69$, Cronach $10^{\mathrm u}$ $28^{\mathrm m}$ $3^{\mathrm s}50$.

\par\indent Die Länge von Cronach ergibt sich also aus dem Vergleich mit

\begin{center}
\fontsize{10}{12}\selectfont\setlength{\tabcolsep}{3pt}\renewcommand{\arraystretch}{1.15}\begin{tabular}{ll}
Göttingen & $35^{\mathrm m}$ $55^{\mathrm s}81$ \\
Mannheim & 35 \quad $57.49$ \\
im Mittel & 35 \quad $56.65$ oder in Graden $28^\circ$ $59'$ $10''$ \\
\end{tabular}
\end{center}
\SourceLeaf{359}\par\indent Der Oberst gibt die Polhöhe an mit $=50^\circ$ $14'$ $42''6$.

\par\indent An unserer neuen Sternwarte wird auch diesen Sommer eifrig gearbeitet. Demnächst erwarte ich einen 12-zölligen Multiplikationskreis von \textsc{Reichenbach}. Herr Dr. \textsc{Gerling} geht als Mathematiklehrer an das Lyzeum in Cassel. Er wird seine Stelle im September antreten. Wegen seines Geschicks und seines Eifers wäre es wünschenswert, dass er auch dort noch astronomisch tätig sein kann. Seine Probeschrift über Sternbedeckungen und Finsternisse und besonders über die bei uns (aber nicht in Berlin) ringförmige Sonnenfinsternis von 1820 (7. Sept.) wird gerade gedruckt.
